\section{Experiments}
\label{sec:experiments}

We ask four questions. (Q1)~Does unrestricted high-rate contact
selection improve success as terrain becomes harder? (Q2)~How much of
the benefit comes from emergent concurrency? (Q3)~Does exhaustive
enumeration matter relative to sampling? (Q4)~Can the frozen planner
run in real time with a model-based stack on a physical robot?

\subsection{Training and Deployment Setup}
\label{subsec:exp-setup}

All conditions use one PPO checkpoint trained as in
\autoref{subsec:method-training}, frozen after training. Conditions
differ only in the action-selection constraints listed below, applied
to the same logits. \pending{Random seeds, the checkpoint hash, and all
launch parameters are logged for every trial.}

\subsection{Gazebo Terrain Benchmark}
\label{subsec:exp-gazebo}

\textit{Terrain.} Each world is a grid of rigid tiles on pillars with
a fraction of tiles removed; removed tiles are true drops rather than
painted regions. Terrain difficulty is the percentage of removed tile
area, from $\corlDifficultyMin\%$ to $\corlDifficultyMax\%$ in
$\corlDifficultyStep\%$ steps. \pending{Tile size, track length, and
the world seeds per difficulty level will be reported here.}

\textit{Protocol.} Commanded forward speed ranges from
$\corlVelocityMin$ to $\corlVelocityMax$\,m/s. Each (method,
difficulty, speed) cell is run for \pending{\corlSimTrialsPerCell}
trials of \pending{\corlSimEpisodeDuration}\,s on matched worlds and
seeds. A trial succeeds if the robot reaches the end of the track
without the base contacting the ground, the controller safety check
triggering, or the base height falling below a fixed threshold.

\textit{Conditions.}
\begin{enumerate}
  \item \textbf{Scheduled NMPC:} the native periodic gait of the
    execution stack with its heuristic foothold projected onto convex
    terrain regions \cite{leggedperceptive}.
  \item \textbf{GaitNet-RoundRobin:} GaitNet footholds and durations,
    but legs are served in a fixed cyclic order.
  \item \textbf{GaitNet-Single:} at most one swinging leg. This matches
    the concurrency of ContactNet's walking configuration
    \cite{bratta_contactnet_2024}; ContactNet has no public
    implementation, so we do not compare against it directly.
  \item \textbf{GaitNet-Diagonal:} a second swing is allowed only for
    the diagonal partner of the leg in the air, a trot-like
    restriction.
  \item \textbf{GaitNet-Full:} the proposed planner, with up to
    \numberstringnum{\corlMaxConcurrentSwingLegs} swinging legs and no
    restriction on which legs overlap.
\end{enumerate}

\begin{table}[t]
  \centering
  \caption{Gazebo success rate, mean over speeds.}
  \label{tab:gazebo}
  \small
  \begin{tabular}{lcc}
    \toprule
    \textbf{Method} & \textbf{All difficulties} &
    \textbf{$\ge\corlHighDifficultyRangeLow\%$} \\
    \midrule
    Scheduled NMPC \cite{leggedperceptive} &
      \pending{\corlGzSchedSuccess\%} & \pending{\corlGzSchedSuccessHard\%} \\
    GaitNet-RoundRobin & \pending{XX\%} & \pending{XX\%} \\
    GaitNet-Single & \pending{\corlGzSingleSuccess\%} & \pending{XX\%} \\
    GaitNet-Diagonal & \pending{\corlGzDiagSuccess\%} & \pending{XX\%} \\
    \textbf{GaitNet-Full} & \pending{\textbf{\corlGzFullSuccess\%}} &
      \pending{\textbf{\corlGzFullSuccessHard\%}} \\
    \bottomrule
  \end{tabular}
\end{table}

\begin{figure}[t]
  \centering
  \pending{%
  \setlength{\fboxsep}{0pt}%
  \fbox{\begin{minipage}[c][0.55\linewidth][c]{0.95\linewidth}
    \centering
    \textit{Success rate vs.\ terrain difficulty (pending)}
  \end{minipage}}%
  }
  \caption{Gazebo success rate versus terrain difficulty for all
    conditions, with 95\% Wilson intervals.}
  \label{fig:gazebo-difficulty}
\end{figure}

\pending{\autoref{tab:gazebo} and \autoref{fig:gazebo-difficulty}
report the results. Describe the ordering of methods, where the gap
opens as difficulty increases, and whether differences are
statistically significant (report the test used).}

\subsection{Concurrency Ablation}
\label{subsec:exp-concurrency}

GaitNet-Single, GaitNet-Diagonal, and GaitNet-Full share the same
weights and differ only in which second swing is admissible, so their
difference isolates the contribution of emergent overlap (Q2). We
report the fraction of time with two legs in swing, the distribution
of overlapping leg pairs, and the number of step requests blocked by
the restriction. \pending{In GaitNet-Full, two legs are in swing
\corlGzOverlapPct\% of the time. Report the pair distribution and
whether same-side pairs occur.} An example swing schedule is shown in
\autoref{fig:swing-schedule}.

\begin{figure}[t]
  \centering
  \pending{%
  \setlength{\fboxsep}{0pt}%
  \fbox{\begin{minipage}[c][0.45\linewidth][c]{0.95\linewidth}
    \centering
    \textit{Gazebo swing schedule (pending)}
  \end{minipage}}%
  }
  \caption{Swing schedule from a GaitNet-Full Gazebo trial. Each row
    is one leg; bars are swings.}
  \label{fig:swing-schedule}
\end{figure}

\subsection{Candidate-Enumeration Study}
\label{subsec:exp-enumeration}

To separate the benefit of dense selection from that of the learned
value function (Q3), we evaluate GaitNet-Full with (i)~exhaustive
enumeration of all valid cells and (ii)~the
\corlTrainTotalCandidates-candidate sampled set used in training, and
\pending{(iii)~intermediate per-leg sample counts}. The weights are
identical across settings. \pending{Dense enumeration reaches
\corlGzDenseSuccess\% success versus \corlGzSampledSuccess\% with
sampling.}

\subsection{Computation and Real-Time Feasibility}
\label{subsec:exp-compute}

We measure (i)~the network forward time, (ii)~the end-to-end service
latency from request to response, and (iii)~the realized decision
interval, which also includes request gating in the controller. In an
isolated process, dense scoring takes \corlDenseScoreMs\,ms.
\pending{In Gazebo, the end-to-end service latency is
\corlGzServiceLatencyMs\,ms (p90 \corlGzServiceLatencyPNinetyMs\,ms),
and the realized decision interval is XX\,ms against the
\corlControlTickMs\,ms target.} The NMPC and WBC run at the rates in
\autoref{tab:stack} regardless of planner latency, since the planner
only updates the mode schedule and foothold targets.

\subsection{Physical-Robot Proof-of-Concept}
\label{subsec:exp-hardware}

We deploy the frozen policy on a Unitree Go1 through the same NMPC/WBC
stack (\autoref{tab:stack}), with an Intel RealSense D435i providing
the point cloud. The track is \verify{\corlHwTerrainBlockMinMm--%
\corlHwTerrainBlockMaxMm\,mm blocks on a mat}, the robot is commanded
\verify{\corlHwTrackLengthM\,m at \corlHwCommandSpeed\,m/s}, and a trial
succeeds if the robot reaches the target distance.

\textit{Deployment adaptations.} The hardware run differs from the
Gazebo configuration in ways that we report explicitly:
\begin{itemize}
  \item Requests are gated to at most one every
    \corlHwRequestPeriod\,s, with a \corlHwNoopCooldown\,s cooldown after
    a no-action response, so the realized decision rate is at most
    4\,Hz rather than \corlControlHz\,Hz.
  \item The requested duration is floored at
    $\max(\corlHwMinSwingDuration,\, d/\corlHwMaxFootSpeed)$\,s, where
    $d$ is the swing travel in meters, and capped at
    \corlHwMaxSwingDuration\,s, because the shortest learned swings
    exceed the Go1's achievable foot speed.
  \item A second swing is restricted to the diagonal partner.
  \item The no-action logit is reduced by \corlHwNoopLogitPenalty{}
    and recently used legs receive a penalty of
    \corlHwRepeatPenalty{}.
\end{itemize}
The policy weights are unchanged; these adaptations act on the
interface, not on the network.

\textit{Results.} \verify{GaitNet completed
\corlHwTerrainSuccesses/\corlHwTerrainAttempts{} attempts on the
block track. The single failure was a step that landed short onto a
block.} \verify{In a trial with diagonal overlap enabled, two legs
were simultaneously in swing \corlHwOverlapPct\% of the time
(\corlHwOverlapPairs{} overlapping pairs, all diagonal; median overlap
\corlHwOverlapMedianMs\,ms, maximum \corlHwOverlapMaxMs\,ms, against
\corlHwNominalSwingMs\,ms swings).} Dense-mode service latency was
\verify{\corlHwServiceLatencyMinMs--\corlHwServiceLatencyMaxMs\,ms}.
\pending{Report the planner latency distribution and mean traversal
time if additional trials are collected.} With $n=\corlHwTerrainAttempts$
we make no statistical comparison against any baseline; these trials
establish only that the frozen policy and interface execute on
physical hardware.

\subsection{Limitations}
\label{subsec:exp-limitations}

\textit{Decision rate on hardware.} The high-rate mechanism evaluated
in Gazebo was not reproduced on the Go1, where request gating limited
decisions to at most 4\,Hz and overlap was brief. The hardware trials
therefore demonstrate feasibility of the interface, not the benefit of
high-rate selection.

\textit{Interface adaptations.} The duration floor, diagonal
restriction, and logit penalties are hand-tuned for the Go1. They
reduce the degrees of freedom that GaitNet exercises on hardware, and
their individual effects have not been ablated.

\textit{Controller dependence.} The planner does not access the
controller's internal optimization, but its success depends on the
controller tracking the requested footholds and durations; we have
evaluated one training controller and one execution stack.

\textit{Concurrency bound.} The rule that
\numberstringnum{\corlMinContactLegs} legs remain in contact is fixed
in all experiments and has not been ablated.

\textit{Perception.} The validity mask is thresholded from a planar
decomposition of the elevation map. Small voids below the camera's
effective resolution are not represented, and the sim-to-real gap in
the mask has not been quantified.
